% Purpose: Complete mathematical, statistical, and graph-theoretic exposition of the pipeline evaluation framework.
% Status: Draft; author and venue review pending.
% Source-of-truth: OGS/src/ogscatalog.py, OGS/src/ogsutils.py, OGS/src/ogsqc.py, OGS/src/ogsconstants.py

\label{chap:methodology}

This chapter establishes the theoretical, statistical, and graph-theoretic evaluation framework used to compare automated machine-learning catalogs with an analyst-curated OGS reference bulletin (\cite{SMINO}). The reference bulletin supplies reviewed labels within its documented scope; it is not evidence that an unmatched automated candidate is physically false.

Evaluating machine-learning models in observational seismology presents a fundamental mathematical departure from standard classification tasks in computer vision or natural language processing:
\begin{enumerate}
  \item \textbf{Continuous Spatiotemporal Continuum}: Earthquakes and seismic phase arrivals are point processes unfolding across a continuous time-space domain $\mathbb{R}^3 \times \mathbb{R}^+$, rather than discrete items drawn from pre-segmented finite sets.
  \item \textbf{Unobservable True Negatives}: In a continuous time series, the absence of an earthquake or phase pick constitutes an uncountably infinite or ill-defined null state ($\epsilon_{NN}$), precluding the standard application of specificity, ROC curves, or total sample accuracy.
  \item \textbf{Multi-Detection Combinatorial Ambiguities}: High-sensitivity deep learning pickers trigger multiple candidate arrivals in close temporal proximity to a single true arrival, creating combinatorial assignment conflicts that require global optimization.
\end{enumerate}

To resolve these challenges, we formulate a bipartite graph matching optimization framework grounded in multi-criteria physical similarity metrics, coupled with a domain-adapted seismic confusion matrix and rigorous residual distribution theory.

% =============================================================================
% SECTION 1: PHYSICAL AND MATHEMATICAL DISTANCE METRICS
% =============================================================================
\section{Physical Distance and Similarity Metrics}
\label{sec:distance_metrics}

Let $\mathcal{B} = \{B_1, B_2, \dots, B_N\}$ denote the set of reference ground-truth records (Base catalog compiled by expert seismic analysts) and $\mathcal{T} = \{T_1, T_2, \dots, T_M\}$ denote the set of candidate detections generated by the automated machine-learning pipeline (Target catalog).

Establishing a valid scientific correspondence between a target detection $T_j$ and a reference entry $B_i$ requires explicit, physically bounded distance metrics evaluated across temporal, spatial, and probabilistic feature spaces.

\subsection{Temporal Discrepancy Metric ($\delta t$)}
Arrival timing fidelity is the primary physical constraint in seismic body-wave propagation. The absolute temporal offset between target and reference timestamps is defined as:
\begin{equation}
  \Delta t(B_i, T_j) = t(T_j) - t(B_i), \quad |\Delta t(B_i, T_j)| = |t(T_j) - t(B_i)|,
\end{equation}
where $t(\cdot)$ represents UTC epoch time in seconds with sub-millisecond precision.

To map this continuous offset into a normalized similarity score bounded in $[0, 1]$, we introduce the linear tolerance function:
\begin{equation}
  \delta t(B_i, T_j) = \begin{cases}
    1 - \dfrac{|t(T_j) - t(B_i)|}{\omega_t}, & \text{if } |t(T_j) - t(B_i)| \le \omega_t, \\[8pt]
    0, & \text{otherwise},
  \end{cases}
\end{equation}
where $\omega_t \in \mathbb{R}^+$ defines the maximum admissible physical tolerance window. In this work:
\begin{itemize}
  \item For phase picks: $\omega_t = 0.5\,\mathrm{s}$ (\texttt{PICK\_TIME\_OFFSET}), corresponding to 50 samples at $100\,\mathrm{Hz}$ and capturing typical onset reading uncertainties.
  \item For seismic events: $\omega_t = 2.0\,\mathrm{s}$ (\texttt{EVENT\_TIME\_OFFSET}), accommodating origin time shifts induced by 1D velocity model simplifications and network geometry variations.
\end{itemize}

\subsection{Geodetic Epicentral and Hypocentral Distances ($\delta h, \delta z, \delta d$)}
Earthquake hypocenters are specified by geodetic coordinates: latitude $\phi$, longitude $\lambda$, and focal depth $z$ below the WGS84 reference ellipsoid.

\subsubsection{Epicentral Geodesic Distance ($\delta h$)}
Planar approximations (e.g. flat-Earth Euclidean metrics) introduce non-negligible geometric distortion over regional network apertures ($> 100\,\mathrm{km}$). Epicentral separation $\delta h(B_i, T_j)$ is computed via Vincenty's inverse geodetic formulae over the WGS84 reference ellipsoid (semi-major axis $a = 6{,}378{,}137.0\,\mathrm{m}$, flattening $f = 1/298.257223563$), as implemented in ObsPy's \texttt{gps2dist\_azimuth}:
\begin{equation}
  \delta h(B_i, T_j) = \frac{1}{1000} \cdot \mathrm{geodesic}_{\mathrm{WGS84}}\left( \phi(B_i), \lambda(B_i), \phi(T_j), \lambda(T_j) \right) \quad [\mathrm{km}].
\end{equation}

\subsubsection{Focal Depth Discrepancy ($\delta z$)}
Focal depth resolution is inherently weaker than horizontal resolution in regional seismology due to station placement on the Earth's surface. Depth discrepancy is defined as:
\begin{equation}
  \delta z(B_i, T_j) = |z(T_j) - z(B_i)| \quad [\mathrm{km}].
\end{equation}

\subsubsection{Hypocentral 3D Spatial Distance ($\delta d$)}
The complete three-dimensional hypocentral Euclidean distance is given by the $L_2$-norm combining horizontal geodetic separation and vertical offset:
\begin{equation}
  \delta d_{3D}(B_i, T_j) = \sqrt{\left( \delta h(B_i, T_j) \right)^2 + \left( \delta z(B_i, T_j) \right)^2} \quad [\mathrm{km}].
\end{equation}

For horizontal event matching, the normalized spatial similarity score is formulated as:
\begin{equation}
  \delta d(B_i, T_j) = \begin{cases}
    1 - \dfrac{\delta h(B_i, T_j)}{\omega_d}, & \text{if } \delta h(B_i, T_j) \le \omega_d, \\[8pt]
    0, & \text{otherwise},
  \end{cases}
\end{equation}
where $\omega_d = 8.0\,\mathrm{km}$ (\texttt{EVENT\_DIST\_OFFSET}) defines the spatial acceptance threshold, reflecting the lateral resolution of regional station networks in the Southeastern Alps.

\subsection{Phase Consistency ($\delta p$) and Confidence Ratio ($\delta \rho$)}
At the individual pick level, two additional dimensionless similarity indicators are evaluated:
\begin{enumerate}
  \item \textbf{Phase Type Indicator ($\delta p$)}: An exact Kronecker delta ensuring physical consistency between wave types:
  \begin{equation}
    \delta p(B_i, T_j) = \mathbb{I}\left( p(B_i) = p(T_j) \right) = \begin{cases}
      1, & \text{if } p(B_i) = p(T_j), \\
      0, & \text{if } p(B_i) \neq p(T_j),
    \end{cases}
  \end{equation}
  where $p(\cdot) \in \{\mathrm{P}, \mathrm{S}\}$.
  \item \textbf{Pick Certainty Ratio ($\delta \rho$)}: Quantifies agreement between the model prediction confidence $\rho(T_j) \in (0, 1]$ and analyst weight certainty $\rho(B_i) \in (0, 1]$:
  \begin{equation}
    \delta \rho(B_i, T_j) = \min\left( 1.0, \frac{\rho(T_j)}{\rho(B_i)} \right).
  \end{equation}
\end{enumerate}

% =============================================================================
% SECTION 2: DOMAIN-SPECIFIC SEISMIC CONFUSION MATRIX
% =============================================================================
\section{Domain-Specific Seismic Confusion Matrix}
\label{sec:confusion_matrix}

In conventional machine learning classification, confusion matrices contrast predicted discrete classes against true labels under the assumption that all samples originate from a fixed, enumerated testing corpus. In continuous seismic time-series processing, this assumption fails.

We establish a domain-specific nomenclature that reformulates classification outcomes to reflect the operational realities of seismic network monitoring:

\begin{itemize}
  \item \textbf{Matched ($\mathrm{MH}$)}: Replaces \textit{True Positive} ($\mathrm{TP}$). A detected pick or event $T_j$ that satisfies all spatiotemporal, station, and phase criteria with respect to an authoritative reference entry $B_i$, and is assigned to $B_i$ under a globally optimal one-to-one mapping:
  \begin{equation}
    B_i \Longleftrightarrow T_j.
  \end{equation}
  \item \textbf{Missed ($\mathrm{MS}$)}: Replaces \textit{False Negative} ($\mathrm{FN}$). An authoritative catalog entry $B_i$ present in the human-curated bulletin that the automated pipeline failed to detect or associate within the admissible tolerance bounds.
  \item \textbf{Proposed ($\mathrm{PS}$)}: Replaces \textit{False Positive} ($\mathrm{FP}$). An automated pick or event detection $T_j$ for which no corresponding counterpart exists in the reference bulletin. In seismic monitoring, proposed events carry dual physical interpretations:
  \begin{enumerate}
    \item Spurious triggers caused by cultural noise, electronic glitches, or teleseismic wave trains.
    \item Legitimate micro-earthquakes ($M_L < 1.0$) previously uncataloged by human analysts due to routine operational bandwidth constraints or visual masking in high-noise environments.
  \end{enumerate}
  \item \textbf{Phase Misclassification ($\epsilon_{\mathrm{PS}}, \epsilon_{\mathrm{SP}}$)}: Off-diagonal pick assignment errors where an arrival is detected at the correct station and time, but misclassified in wave type:
  \begin{itemize}
    \item $\epsilon_{\mathrm{PS}}$: A true P-wave in the reference catalog misidentified as an S-wave by the neural network.
    \item $\epsilon_{\mathrm{SP}}$: A true S-wave in the reference catalog misidentified as a P-wave by the neural network.
  \end{itemize}
  \item \textbf{Unobservable Null Set ($\epsilon_{\mathrm{NN}}$)}: Replaces \textit{True Negative} ($\mathrm{TN}$). Represents intervals of continuous waveform data where both the reference catalog and the automated model agree that no earthquake or phase arrival occurred. Because a continuous time continuum contains infinitely many potential non-arrival instances, $\epsilon_{\mathrm{NN}}$ cannot be objectively counted:
  \begin{equation}
    \epsilon_{\mathrm{NN}} \equiv \emptyset \quad (\text{Mathematically uncountably infinite / undefined}).
  \end{equation}
  Consequently, all performance metrics dependent on $\mathrm{TN}$ (such as accuracy or specificity) are mathematically invalid in this context.
\end{itemize}

\subsection{Dual Confusion Matrix Formulations}
Tables~\ref{tab:pick_confusion_matrix} and \ref{tab:event_confusion_matrix} define the formal mathematical confusion matrix architectures for phase picks and seismic events, respectively.

\begin{table}[htbp]
  \centering
  \small
  \caption{Domain-specific 3$\times$3 confusion matrix architecture for seismic phase picks.}
  \label{tab:pick_confusion_matrix}
  \begin{threeparttable}
    \begin{tabularx}{\linewidth}{l c c c X}
      \toprule
      & \multicolumn{3}{c}{\textbf{Automated ML Pipeline (Target $T_j$)}} & \\
      \cmidrule(lr){2-4}
      \textbf{OGS Bulletin (Base $B_i$)\tnote{a}} & \textbf{P Pick} & \textbf{S Pick} & \textbf{None (Missed)} & \textbf{Class Totals} \\
      \midrule
      \textbf{P Pick} & $\mathrm{MH}_{\mathrm{P}}$ & $\epsilon_{\mathrm{PS}}$ & $\mathrm{MS}_{\mathrm{P}}$ & $N_{\mathrm{Base}, \mathrm{P}} = \mathrm{MH}_{\mathrm{P}} + \epsilon_{\mathrm{PS}} + \mathrm{MS}_{\mathrm{P}}$ \\
      \textbf{S Pick} & $\epsilon_{\mathrm{SP}}$ & $\mathrm{MH}_{\mathrm{S}}$ & $\mathrm{MS}_{\mathrm{S}}$ & $N_{\mathrm{Base}, \mathrm{S}} = \epsilon_{\mathrm{SP}} + \mathrm{MH}_{\mathrm{S}} + \mathrm{MS}_{\mathrm{S}}$ \\
      \textbf{None (Proposed)\tnote{b}} & $\mathrm{PS}_{\mathrm{P}}$ & $\mathrm{PS}_{\mathrm{S}}$ & $[\epsilon_{\mathrm{NN}}]$\tnote{c} & $N_{\mathrm{Target}, \mathrm{Proposed}} = \mathrm{PS}_{\mathrm{P}} + \mathrm{PS}_{\mathrm{S}}$ \\
      \midrule
      \textbf{Target Totals} & $N_{\mathrm{Tgt}, \mathrm{P}} = \mathrm{MH}_{\mathrm{P}} + \epsilon_{\mathrm{SP}} + \mathrm{PS}_{\mathrm{P}}$ & $N_{\mathrm{Tgt}, \mathrm{S}} = \epsilon_{\mathrm{PS}} + \mathrm{MH}_{\mathrm{S}} + \mathrm{PS}_{\mathrm{S}}$ & --- & $N_{\mathrm{Total}} = \sum \mathrm{MH} + \sum \mathrm{MS} + \sum \mathrm{PS}$ \\
      \bottomrule
    \end{tabularx}
    \begin{tablenotes}[flushleft]
      \footnotesize
      \item[a] [Analyst-Normalized] Human-curated ground truth from the authoritative OGS seismic bulletin.
      \item[b] [Model Forecast / Hypothesis] Algorithmic phase picks proposed by the deep learning picker with no bulletin counterpart.
      \item[c] $[\epsilon_{\mathrm{NN}}]$ denotes the unobservable continuous time background; excluded from quantitative metric denominators.
      \item Target \& Scope: Station-by-station phase pick classification evaluated across the continuous multi-year test archive.
      \item Notice how off-diagonal cells $\epsilon_{\mathrm{PS}}$ and $\epsilon_{\mathrm{SP}}$ capture physical wave-type misidentification.
      \item Evidence Anchor: Implemented in \texttt{OGSCatalog.bgmaPicks} within \texttt{OGS/src/ogscatalog.py}.
    \end{tablenotes}
  \end{threeparttable}
\end{table}

\begin{table}[htbp]
  \centering
  \small
  \caption{Domain-specific 2$\times$2 confusion matrix architecture for seismic events.}
  \label{tab:event_confusion_matrix}
  \begin{threeparttable}
    \begin{tabularx}{\linewidth}{l c c X}
      \toprule
      & \multicolumn{2}{c}{\textbf{Automated Pipeline (Target $T_j$)}} & \\
      \cmidrule(lr){2-3}
      \textbf{OGS Bulletin (Base $B_i$)\tnote{a}} & \textbf{Associated Event} & \textbf{None (Missed)} & \textbf{Class Totals} \\
      \midrule
      \textbf{Associated Event} & $\mathrm{MH}_{\mathrm{E}}$ & $\mathrm{MS}_{\mathrm{E}}$ & $N_{\mathrm{Base}, \mathrm{Events}} = \mathrm{MH}_{\mathrm{E}} + \mathrm{MS}_{\mathrm{E}}$ \\
      \textbf{None (Proposed)\tnote{b}} & $\mathrm{PS}_{\mathrm{E}}$ & $[\epsilon_{\mathrm{NN}}]$\tnote{c} & $N_{\mathrm{Proposed}, \mathrm{Events}} = \mathrm{PS}_{\mathrm{E}}$ \\
      \midrule
      \textbf{Target Totals} & $N_{\mathrm{Tgt}, \mathrm{Events}} = \mathrm{MH}_{\mathrm{E}} + \mathrm{PS}_{\mathrm{E}}$ & --- & $N_{\mathrm{Total\,Events}} = \mathrm{MH}_{\mathrm{E}} + \mathrm{MS}_{\mathrm{E}} + \mathrm{PS}_{\mathrm{E}}$ \\
      \bottomrule
    \end{tabularx}
    \begin{tablenotes}[flushleft]
      \footnotesize
      \item[a] [Analyst-Normalized] Verified seismic events cataloged in the OGS reference bulletin.
      \item[b] [Model Forecast / Hypothesis] Automatically formed event clusters with preliminary or NonLinLoc hypocenters.
      \item[c] Unobservable non-event space; excluded from statistical evaluations.
      \item Target \& Scope: Network-level earthquake association and catalog recovery.
      \item Notice that event matching requires multi-station pick coherence, making $\mathrm{PS}_{\mathrm{E}}$ far less noisy than pick-level $\mathrm{PS}$.
      \item Evidence Anchor: Implemented in \texttt{OGSCatalog.bgmaEvents} within \texttt{OGS/src/ogscatalog.py}.
    \end{tablenotes}
  \end{threeparttable}
\end{table}

% =============================================================================
% SECTION 3: BIPARTITE GRAPH MATCHING THEORY
% =============================================================================
\section{Bipartite Graph Matching Theory}
\label{sec:bipartite_matching}

When an automated model processes dense seismic sequences, multiple candidate picks or associated events frequently fall within the tolerance window of a single reference entry. A naive greedy matching strategy---such as pairing each catalog entry with its nearest temporal neighbor---produces suboptimal cascading assignments: an early greedy match can consume a candidate that would have provided the unique valid match for an adjacent record.

To guarantee globally optimal, non-conflicting, one-to-one correspondences, we model the benchmark comparison as a \textbf{Maximum-Weight Bipartite Graph Matching} optimization problem.

\subsection{Formal Graph Formulation}
Let $G = (V, E, W)$ define an undirected, weighted bipartite graph whose node set $V$ is partitioned into two disjoint subsets:
\begin{equation}
  V = V_B \cup V_T, \quad V_B \cap V_T = \emptyset,
\end{equation}
where $V_B = \{v_1^B, \dots, v_N^B\}$ represents the $N$ reference records in the Base catalog and $V_T = \{v_1^T, \dots, v_M^T\}$ represents the $M$ predictions in the Target catalog.

An undirected edge $e_{ij} = (v_i^B, v_j^T) \in E$ exists between base node $v_i^B$ and target node $v_j^T$ if and only if the candidate pair satisfies all strict domain admissibility constraints:
\begin{enumerate}
  \item \textbf{Pick Matching Feasibility}:
  \begin{equation}
    e_{ij} \in E_{\mathrm{Picks}} \iff \begin{cases}
      \mathrm{station}(B_i) = \mathrm{station}(T_j), \\
      |t(T_j) - t(B_i)| \le \omega_t \quad (\omega_t = 0.5\,\mathrm{s}).
    \end{cases}
  \end{equation}
  \item \textbf{Event Matching Feasibility}:
  \begin{equation}
    e_{ij} \in E_{\mathrm{Events}} \iff \begin{cases}
      |t(T_j) - t(B_i)| \le \omega_t \quad (\omega_t = 2.0\,\mathrm{s}), \\
      \delta h(B_i, T_j) \le \omega_d \quad (\omega_d = 8.0\,\mathrm{km}).
    \end{cases}
  \end{equation}
\end{enumerate}

\subsection{Composite Edge Weight Scoring Functions}
Every admissible edge $e_{ij} \in E$ is assigned a scalar weight $W(e_{ij}) \in (0, 1]$ representing the composite spatiotemporal similarity between records $B_i$ and $T_j$.

\subsubsection{Pick-Level Similarity Function ($\Delta_{\mathrm{Picks}}$)}
For phase picks, the composite similarity score combines temporal proximity (97\%), phase agreement (2\%), and pick certainty (1\%), as implemented in \texttt{OGS/src/ogsutils.py} (lines 452--472):
\begin{equation}
  \Delta_{\mathrm{Picks}}(B_i, T_j) = w_t \cdot \delta t(B_i, T_j) + w_p \cdot \delta p(B_i, T_j) + w_\rho \cdot \delta \rho(B_i, T_j),
\end{equation}
subject to the normalization constraint:
\begin{equation}
  w_t + w_p + w_\rho = 1.0, \quad w_t = 0.97, \quad w_p = 0.02, \quad w_\rho = 0.01.
\end{equation}
The overwhelming weight assigned to $w_t$ ensures that timing precision dominates assignment decisions, while the small positive weight $w_p$ breaks ties in favor of matching identical phase labels.

\subsubsection{Event-Level Similarity Function ($\Delta_{\mathrm{Events}}$)}
For seismic events, the composite similarity score balances origin time consistency (99\%) and geodetic epicentral proximity (1\%), as implemented in \texttt{OGS/src/ogsutils.py} (lines 475--495):
\begin{equation}
  \Delta_{\mathrm{Events}}(B_i, T_j) = w_t \cdot \delta t(B_i, T_j) + w_d \cdot \delta d(B_i, T_j),
\end{equation}
with normalized coefficients:
\begin{equation}
  w_t + w_d = 1.0, \quad w_t = 0.99, \quad w_d = 0.01.
\end{equation}

\subsection{Maximum Weight Matching Optimization}
A matching $M^* \subseteq E$ is an edge subset such that no two edges share a common vertex:
\begin{equation}
  \forall e_1 = (u_1, v_1), \, e_2 = (u_2, v_2) \in M^*, \quad e_1 \neq e_2 \implies \{u_1, v_1\} \cap \{u_2, v_2\} = \emptyset.
\end{equation}

The optimal one-to-one correspondence is formulated as finding the matching $M^*$ that maximizes the total sum of edge weights:
\begin{equation}
  M^* = \arg\max_{M \subseteq E} \sum_{e \in M} W(e).
\end{equation}

This optimization is solved via the primal-dual Hungarian (Kuhn-Munkres) algorithm or Galil's $O(|V| \cdot |E| \log |V|)$ maximum weight matching algorithm, implemented via NetworkX's \texttt{max\_weight\_matching(G, maxcardinality=False, weight='weight')}. The condition \texttt{maxcardinality=False} is vital: it allows the algorithm to leave poor, marginal edges unmatched if assigning them would reduce overall weight, preserving high fidelity.

Following optimization, the node partitions are resolved deterministically:
\begin{align}
  \mathcal{M} &= \left\{ (B_i, T_j) \mid (v_i^B, v_j^T) \in M^* \right\} \quad &(\textbf{Matched Set, MH}), \\
  \mathcal{MS} &= \left\{ B_i \in \mathcal{B} \mid \forall v_j^T \in V_T, \, (v_i^B, v_j^T) \notin M^* \right\} \quad &(\textbf{Missed Set, MS}), \\
  \mathcal{PS} &= \left\{ T_j \in \mathcal{T} \mid \forall v_i^B \in V_B, \, (v_i^B, v_j^T) \notin M^* \right\} \quad &(\textbf{Proposed Set, PS}).
\end{align}

\subsection{Synthetic Worked Example and Eligibility Accounting}
\label{subsec:bgma_worked_example}
The following deliberately synthetic example mirrors the station-and-time eligibility behavior tested in \texttt{OGS/test/testogscatalog.py}; it is not a seismic result. Let the analyst reference catalog contain a P pick at station \texttt{AAA} at $t_0=$ 00:00:00 and an S pick at the same station at $t_1=$ 00:00:10. Let the automated catalog contain a P pick at \texttt{AAA} at $t_0+0.2\,\unit{\second}$, an S pick at \texttt{AAA} at $t_1+0.3\,\unit{\second}$, a second S candidate at $t_1+0.8\,\unit{\second}$, and a P candidate at unlisted station \texttt{BBB} at $t_0+0.1\,\unit{\second}$. With $\omega_t=0.5\,\unit{\second}$ and an inventory containing only \texttt{AAA}, the first two candidates form the two admissible edges and are matched. The late S candidate has no feasible edge and enters $\mathrm{PS}$; the \texttt{BBB} candidate is excluded before matching and enters $\mathrm{PicksSP}$.

Thus the BGMA-eligible counts are $\mathrm{MH}=2$, $\mathrm{MS}=0$, $\mathrm{PS}=1$, and $\mathrm{PicksSP}=1$. The eligible-reference recall is $2/(2+0)=1$, while the unmatched-reference rate among eligible target candidates is $1/(2+1)=1/3$. Reporting $\mathrm{PicksSP}$ separately is essential: it is neither a match nor a proposed pick, and must not enter either denominator. The production implementation analogously records spatially excluded events in \texttt{EventsSM} and \texttt{EventsSP}.

% =============================================================================
% SECTION 4: RESIDUAL DISTRIBUTION ANALYSIS
% =============================================================================
\section{Residual Distribution Analysis}
\label{sec:residuals}

For all matched pairs $(B_i, T_j) \in \mathcal{M}$, physical discrepancies are evaluated across temporal, spatial, and magnitude dimensions to detect systematic instrumentation biases, velocity model deficiencies, or site-response errors.

\subsection{Temporal Residuals ($\Delta t$)}
For matched picks and origin times, the signed temporal residual is defined as:
\begin{equation}
  \Delta t_k = t(T_k) - t(B_k), \quad \forall k \in \{1, \dots, |\mathcal{M}|\}.
\end{equation}
A positive residual ($\Delta t > 0$) indicates that the automated model predicts an arrival later than the human analyst, whereas a negative residual ($\Delta t < 0$) indicates an earlier trigger.

To characterize the residual distribution, we compute four complementary statistical moments and robust metrics:
\begin{itemize}
  \item \textbf{Mean Error (Bias)}:
  \begin{equation}
    \mu_{\Delta t} = \frac{1}{|\mathcal{M}|} \sum_{k=1}^{|\mathcal{M}|} \Delta t_k.
  \end{equation}
  A non-zero mean reveals systematic delay, such as causal filtering latency or systematic phase identification late on the emergent wave train.
  \item \textbf{Standard Deviation (Timing Jitter)}:
  \begin{equation}
    \sigma_{\Delta t} = \sqrt{\frac{1}{|\mathcal{M}| - 1} \sum_{k=1}^{|\mathcal{M}|} (\Delta t_k - \mu_{\Delta t})^2}.
  \end{equation}
  \item \textbf{Mean Absolute Error (MAE, Outlier-Robust)}:
  \begin{equation}
    \mathrm{MAE}_{\Delta t} = \frac{1}{|\mathcal{M}|} \sum_{k=1}^{|\mathcal{M}|} |\Delta t_k|.
  \end{equation}
  \item \textbf{Root Mean Square Error (RMSE, Penalty for Extremes)}:
  \begin{equation}
    \mathrm{RMSE}_{\Delta t} = \sqrt{\frac{1}{|\mathcal{M}|} \sum_{k=1}^{|\mathcal{M}|} (\Delta t_k)^2}.
  \end{equation}
\end{itemize}

\subsection{Spatial Epicentral and Hypocentral Residuals ($\Delta h, \Delta z, \Delta d$)}
For matched earthquake events, spatial residuals evaluate the relocation accuracy of the automated pipeline relative to the OGS bulletin:
\begin{align}
  \Delta h_k &= \delta h(B_k, T_k) \quad &(\text{Epicentral Horizontal Residual}), \\
  \Delta z_k &= z(T_k) - z(B_k) \quad &(\text{Signed Vertical Depth Residual}), \\
  \Delta d_k &= \sqrt{(\Delta h_k)^2 + (\Delta z_k)^2} \quad &(\text{3D Hypocentral Separation}).
\end{align}
Correlating spatial residuals with station network azimuthal gap ($\mathrm{GAP}$) and secondary azimuth distributions isolates location biases introduced by asymmetric station geometries along the Alpine arc.

\subsection{Magnitude Residuals ($\Delta M_L$)}
For matched events possessing valid local magnitude estimates in both catalogs, the magnitude residual is defined as:
\begin{equation}
  \Delta M_{L, k} = M_L(T_k) - M_L(B_k).
\end{equation}
Statistical evaluation of $\mu_{\Delta M_L}$ and $\sigma_{\Delta M_L}$ validates whether the synthetic Wood-Anderson simulation and regional attenuation parameters ($\log_{10} A_0$) reproduce the historical OGS magnitude scale without systematic scaling drift.

% =============================================================================
% SECTION 5: MACHINE LEARNING EVALUATION METRICS
% =============================================================================
\section{Machine Learning Performance Metrics}
\label{sec:ml_metrics}

Using the domain-adapted confusion matrix and the partition cardinalities resolved by bipartite matching, we define standardized performance metrics adapted for continuous seismic monitoring.

\subsection{Recall (Sensitivity / Completeness)}
Recall quantifies the pipeline's capability to successfully recover verified ground-truth observations from the continuous stream:
\begin{itemize}
  \item \textbf{P-Wave Pick Recall}:
  \begin{equation}
    \mathrm{Recall}_{\mathrm{P}} = \frac{\mathrm{MH}_{\mathrm{P}}}{\mathrm{MH}_{\mathrm{P}} + \mathrm{MS}_{\mathrm{P}} + \epsilon_{\mathrm{PS}}}.
  \end{equation}
  \item \textbf{S-Wave Pick Recall}:
  \begin{equation}
    \mathrm{Recall}_{\mathrm{S}} = \frac{\mathrm{MH}_{\mathrm{S}}}{\mathrm{MH}_{\mathrm{S}} + \mathrm{MS}_{\mathrm{S}} + \epsilon_{\mathrm{SP}}}.
  \end{equation}
  \item \textbf{Total Pick Recall}:
  \begin{equation}
    \mathrm{Recall}_{\mathrm{Picks}} = \frac{\mathrm{MH}_{\mathrm{P}} + \mathrm{MH}_{\mathrm{S}}}{\left(\mathrm{MH}_{\mathrm{P}} + \mathrm{MS}_{\mathrm{P}} + \epsilon_{\mathrm{PS}}\right) + \left(\mathrm{MH}_{\mathrm{S}} + \mathrm{MS}_{\mathrm{S}} + \epsilon_{\mathrm{SP}}\right)}.
  \end{equation}
  \item \textbf{Event Catalog Recall}:
  \begin{equation}
    \mathrm{Recall}_{\mathrm{Events}} = \frac{\mathrm{MH}_{\mathrm{E}}}{\mathrm{MH}_{\mathrm{E}} + \mathrm{MS}_{\mathrm{E}}}.
  \end{equation}
\end{itemize}

A high event recall ($\ge 90\%$) indicates that the automated pipeline reliably recovers the vast majority of earthquakes identified by human analysts.

\subsection{Precision (Positive Predictive Value)}
Precision measures the fraction of pipeline detections that represent verified bulletin events:
\begin{itemize}
  \item \textbf{Pick-Level Precision}:
  \begin{equation}
    \mathrm{Precision}_{\mathrm{P}} = \frac{\mathrm{MH}_{\mathrm{P}}}{\mathrm{MH}_{\mathrm{P}} + \mathrm{PS}_{\mathrm{P}} + \epsilon_{\mathrm{SP}}}, \quad \mathrm{Precision}_{\mathrm{S}} = \frac{\mathrm{MH}_{\mathrm{S}}}{\mathrm{MH}_{\mathrm{S}} + \mathrm{PS}_{\mathrm{S}} + \epsilon_{\mathrm{PS}}}.
  \end{equation}
  \item \textbf{Event-Level Precision}:
  \begin{equation}
    \mathrm{Precision}_{\mathrm{Events}} = \frac{\mathrm{MH}_{\mathrm{E}}}{\mathrm{MH}_{\mathrm{E}} + \mathrm{PS}_{\mathrm{E}}}.
  \end{equation}
\end{itemize}

\subsection{Unmatched-Reference Rate and Conditional False Discovery Rate}
In operational seismic monitoring centers, minimizing analyst alert fatigue is critical. Before independent adjudication of proposed candidates, the following quantity is an \emph{unmatched-reference rate}: it quantifies the proportion of eligible automated detections that lack a bulletin counterpart, but does not establish that they are false discoveries.
\begin{equation}
  \mathrm{FDR} = 1 - \mathrm{Precision} = \frac{\mathrm{PS}}{\mathrm{MH} + \mathrm{PS}}.
\end{equation}
At the event level:
\begin{equation}
  \mathrm{FDR}_{\mathrm{Events}} = \frac{\mathrm{PS}_{\mathrm{E}}}{\mathrm{MH}_{\mathrm{E}} + \mathrm{PS}_{\mathrm{E}}}.
\end{equation}
  It can be called an FDR only after a documented independent review classifies the relevant proposed candidates as false or true events. Every report of this rate must state the inventory and spatial exclusions ($\mathrm{PicksSM}$, $\mathrm{PicksSP}$, $\mathrm{EventsSM}$, and $\mathrm{EventsSP}$).

\subsection{$F_1$-Score (Harmonic Mean)}
The $F_1$-score provides a balanced harmonic metric integrating sensitivity and positive predictive value:
\begin{equation}
  F_1 = 2 \cdot \frac{\mathrm{Precision} \cdot \mathrm{Recall}}{\mathrm{Precision} + \mathrm{Recall}} = \frac{2 \cdot \mathrm{MH}}{2 \cdot \mathrm{MH} + \mathrm{PS} + \mathrm{MS} + \epsilon_{\mathrm{mismatch}}}.
\end{equation}

\subsection{Operational Sensitivity--Specificity Trade-Off}
In operational practice, varying the model detection threshold $\tau_{\mathrm{threshold}}$ sweeps a continuous curve through metric space:
\begin{itemize}
  \item Lowering $\tau_{\mathrm{threshold}} \to 0.1$ maximizes Recall at the expense of elevated FDR, useful for retrospective research aiming for catalog completeness.
  \item Raising $\tau_{\mathrm{threshold}} \to 0.5$ suppresses FDR to $< 5\%$ at the expense of modest Recall degradation, ideal for real-time civil protection alerting.
\end{itemize}

This evaluation framework provides the mathematical and reproducible foundation upon which all subsequent pipeline benchmarks and cross-model comparisons are conducted.
